\chapter{Acute Abdomen}

\label{ch:acute_abdomen}

The term "acute abdomen" refers to a rapid, sudden onset of severe abdominal pain that often signals a surgical emergency. Correct modality selection and systematic image reading are vital to differentiate self-limiting conditions from life-threatening emergencies (such as bowel perforation, bowel obstruction, or mesenteric ischemia).


\section*{1. Diagnostic Modality Selection}
\begin{itemize}
    \item \textbf{Erect Chest X-Ray}: Can show free intraperitoneal air beneath the diaphragm when the patient can be positioned upright, but a negative radiograph does not exclude perforation. CT is more sensitive for small amounts of free air when clinically appropriate.
    \item \textbf{Abdominal X-Ray (AXR)}: Has limited clinical utility but remains useful for evaluating suspected bowel obstruction, sigmoid/cecal volvulus, or radiopaque foreign bodies.
    \item \textbf{Ultrasound (US)}: Often first-line for suspected biliary disease and selected acute pelvic or pregnancy-related questions. A nondiagnostic examination does not exclude ectopic pregnancy, ovarian torsion, or pelvic inflammatory disease when clinical concern remains.
    \item \textbf{CT Abdomen and Pelvis}: Frequently appropriate for undifferentiated acute abdominal pain in nonpregnant adults, with the contrast protocol tailored to the suspected cause. Ultrasound or MRI may be preferred in pregnancy, children, and specific clinical scenarios.
\end{itemize}


\section*{2. Key Pathological Entities and Radiologic Findings}

\subsection*{Pneumoperitoneum (Bowel Perforation)}
\begin{itemize}
    \item \textbf{Erect CXR}: Look for a thin, dark crescent of air separating the right hemidiaphragm from the liver (crescent sign).
    \item \textbf{CT Abdomen}: Highly sensitive for trace amounts of free air, which pools anteriorly in the supine patient (just beneath the anterior abdominal wall).
\end{itemize}

\subsection*{Bowel Obstruction}
\begin{itemize}
    \item \textbf{Small Bowel Obstruction (SBO)}:
    \item Dilated small bowel loops ($> 3$ cm in diameter).
    \item Prominent \textbf{valvulae conniventes} (folds that stretch completely across the lumen like a ladder).
    \item Presence of multiple air-fluid levels (especially on erect AXR) and a clear \textbf{transition point} on CT between dilated proximal loops and collapsed distal loops.
    \item \textbf{Large Bowel Obstruction (LBO)}:
    \item Dilated colon ($> 6$ cm; cecum $> 9$ cm).
    \item Presence of \textbf{haustra} (folds that do not stretch completely across the lumen).
    \item Sigmoid Volvulus: Presents with a massive, dilated sigmoid loop arising from the pelvis, resembling a \textbf{"coffee bean"} on AXR.
\end{itemize}

\subsection*{Acute Appendicitis}
Look for a blind-ending tubular structure in the right lower quadrant:
\begin{itemize}
    \item Outer diameter \textbf{$> 6$ mm}.
    \item Wall thickening ($> 2$ mm) showing a hyper-enhancing "target sign" on contrast-enhanced CT.
    \item Surrounding appendiceal fat stranding and the presence of a calcified \textbf{appendicolith} at the base.
\end{itemize}

\subsection*{Acute Diverticulitis}
Usually localized to the sigmoid colon:
\begin{itemize}
    \item Focal sigmoid colon wall thickening with multiple diverticula out-pouchings.
    \item Pronounced surrounding mesenteric \textbf{fat stranding} (indicating active inflammation).
    \item Check for complications: local abscess (fluid collection with gas) or extraluminal free air.
\end{itemize}

\subsection*{Acute Pancreatitis}
Diagnosis is usually based on clinical and laboratory criteria. Contrast-enhanced CT is reserved for diagnostic uncertainty, suspected complications, severe disease, or failure to improve as expected; very early CT may underestimate necrosis. When CT is obtained, it can characterize complications and extent:
\begin{itemize}
    \item Pancreatic enlargement with indistinct borders and surrounding fat stranding/fluid collections.
    \item \textbf{Necrotizing Pancreatitis}: Indicated by areas of non-enhancing pancreatic parenchyma on contrast CT.
\end{itemize}

\subsection*{Urolithiasis (Renal Colic)}
\begin{itemize}
    \item \textbf{Non-Contrast CT (NCCT)}: Commonly used for suspected acute ureteral stone in nonpregnant adults. A calculus is identified by its location and attenuation relative to surrounding tissues; there is no universal HU threshold required for diagnosis.
    \item \textbf{Secondary Signs}: Hydroureteronephrosis (dilated collecting system) and perinephric fat stranding.
\end{itemize}


\section*{3. Structured Reporting Example: Small Bowel Obstruction (SBO)}

\begin{testbox}[Indications:]
68-year-old female presenting with severe abdominal distension, obstipation, and bilious vomiting. History of prior open cholecystectomy.
\end{testbox}
\begin{refbox}[Findings:]
\textbf{Solid Organs}: Liver, gallbladder, spleen, pancreas, and adrenal glands are normal. No hydronephrosis. Prior cholecystectomy surgical clips noted in the right upper quadrant.\\
\textbf{Gastrointestinal Tract}: There is marked dilatation of multiple small-bowel loops, measuring up to 4.2 cm, with air-fluid levels. A discrete transition point is present in the right lower quadrant. Given the prior surgery and no other visible cause, an adhesive obstruction is favored; individual adhesive bands are not usually directly visualized. Distal small bowel and colon are decompressed. The bowel wall enhances without CT findings of pneumatosis or portal venous gas. The appendix is normal.\\
\textbf{Vessels \& Peritoneum}: The abdominal aorta is normal. There is a small amount of simple free fluid in the pelvis. No pneumoperitoneum detected.
\end{refbox}
\begin{pearlbox}[Impression:]
1. High-grade mechanical small bowel obstruction with a transition point in the right lower quadrant, likely secondary to post-surgical adhesions.
2. No CT evidence of bowel ischemia or perforation; correlate with examination and laboratory findings if concern persists.
\end{pearlbox}

\section*{Selected Guidance}
\begin{itemize}
    \item American College of Radiology. \href{https://www.radiologyinfo.org/en/info/acs-acute-nonlocalized-abdominal-pain}{Appropriateness guidance for acute nonlocalized abdominal pain}.
    \item American College of Radiology. \href{https://acsearch.acr.org/docs/69474/Narrative/}{Appropriateness Criteria: Right Upper Quadrant Pain}.
    \item American College of Gastroenterology. \href{https://pubmed.ncbi.nlm.nih.gov/38857482/}{2024 guideline on management of acute pancreatitis}.
\end{itemize}

\clearpage
